\documentclass[10pt]{article}

% Specify the margins and text size.
\setlength{\textwidth}{6.4in}
\setlength{\textheight}{9.5in}
\setlength{\oddsidemargin}{0pt}
\setlength{\evensidemargin}{0pt}
\setlength{\topmargin}{0pt}
\setlength{\hoffset}{.05in}
\setlength{\voffset}{-1in}

\setlength{\parskip}{5pt}
\setlength{\parindent}{0pt}

% Load some fonts and symbol packages
\usepackage{latexsym}
\usepackage{pifont}       % contains 'star' symbol for counterinsurgency handbook title
\usepackage{yfonts} 
\usepackage{amsmath}
\usepackage{amsfonts}

\usepackage{graphicx}     % actually, this is loaded by pstricks
\usepackage[T1]{fontenc}
\usepackage{ifthen}
\usepackage{pstricks,pst-grad,pst-text,pst-node,multido,pst-plot,calc,pst-3dplot}
%\usepackage[all]{xy}
%\usepackage{animate}

% The hyperref package inserts links into the pdf.
\definecolor{MyLinkColor}{rgb}{.1,.2,1}
\definecolor{MyCiteColor}{rgb}{.1,1,.2}
\definecolor{MyURLColor}{rgb}{.4,.4,.4}
\usepackage[backref=true,pagebackref=false,hyperindex,colorlinks=true,
  linkcolor=MyLinkColor,urlcolor=MyURLColor]{hyperref}


% The tweaklist package is something I found on the web.  It provides a simple interface
% for making changes to spacing used in the itemize and enumerate environments.  Comment
% this out if you don't care to use tweaklist.
\usepackage{tweaklist}
\renewcommand{\itemhook}{\setlength{\parskip}{2pt}\setlength{\parsep}%
{1pt}\setlength{\topsep}{0pt}\setlength{\itemsep}{0pt}}

\newcommand{\U}{\underline{\hspace{5pt}}}

\usepackage{listings}
\newcommand{\Z}{\hphantom{0}}

\begin{document}
\pagestyle{empty}
\lstset{language=R, showspaces=false, showstringspaces=false}

\href{http://www.shepherd.edu}{\includegraphics[height=1.75cm]{logo-high-res.eps}}
\vspace{-1.69cm}

{\small
\begin{tabular}{cl}
& Math 314\\
& Statistics\\
\hspace{5.28in} & 19 September 2011
\end{tabular}
}
\setlength{\baselineskip}{1.05\baselineskip}
\bigskip

\begin{center}
\textbf{\large  Exam I}
\end{center}
\medskip

1. For the $(x,y)$ data shown in the table, do the following.

\hspace{20pt} (a) Draw a scatter plot. (10 points)

\hspace{20pt} (b) Calculate the correlation coefficient. (10 points)

\hspace{20pt} (c) Find an equation for the SD line. (10 points)

\hspace{20pt} (d) Plot the SD line.  (5 points)

\rput(14.3,2){\begin{tabular}{|r|lllllll|}\hline
$x$ & 1 & 2 & 3 & 4 & 5 & 6 & 7\vphantom{\Large Y}\\[3pt]
$y$ & 2 & 1 & 4 & 3 & 7 & 5 & 6\\[1pt]\hline
\end{tabular}}

\vfill
\eject

2. (Hypothetical)  In a clinical trial, data collection usually starts at
\lq\lq baseline,\rq\rq\ when the subjects are 
recruited into the trial but before they are randomized to 
treatment or control.  Data collection continues until the end of the followup.
Two clinical trials on prevention of heart attacks report baseline data 
on weight, shown below.  In one of these trial, the randomization did not work.
Which  one and why? (15~points)\vspace{-15pt}
\newcommand{\Zc}{\hphantom{0,}}
\begin{center}
\begin{tabular}{cc}
\begin{tabular}{l}
\\
\\
\rput(-12pt,-3pt){(i)}
Treatment\vphantom{\Large Y}\\
Control\\[10pt]
\rput(-12pt,-3pt){(ii)}
Treatment\\
Control
\end{tabular}
&
\begin{tabular}{ccc}
Number of & Average &\\
persons & weight    & SD\\\hline
1,012   & 185 lb    & 25 lb\vphantom{\Large Y}\\
\Zc997  & 143 lb    & 26 lb\\[10pt]
\Zc800  & 166 lb    & 27 lb\\
1,002   & 162 lb    & 25 lb
\end{tabular}
\end{tabular}
\end{center}
\vspace{1.5in}

3. In 2005, men averaged about 540 on the Math SAT.  Estimate the percentage of 
men who scored over 660.  You may assume (i) that the histogram 
followed the normal curve and (ii)~the SD was about 120. (15~points)
\vspace{2.25in}

4. For the list 1, 3, 5, 7, 9, the mean is 5 and the SD is about 2.8.

\hspace{20pt} a) What are the median, minimum, maximum, 25th percentile and the 75th percentile?
(5 points)
\vspace{1in}

\hspace{20pt} b) What are the mean and SD of the list 5, 7, 9, 11, 13?  
Explain briefly. (5 points)
\vspace{0.75in}

\hspace{20pt} c) What are the mean and SD of the list 2, 6, 10, 14, 18?  Explain briefly. (5 points)
\vfill
\eject


5. The table below shows results from two experiments designed to test the Salk vaccine in  1954.
\vspace{-5pt}
\begin{center}
{%\footnotesize
\begin{tabular}{cc}
Randomized controlled          & \\
double-blind experiment        & NFIP study\\
   \begin{tabular}{lcc}\hline
                 & Size      & Rate${}^*$\vphantom{\LARGE Y}\\
      \color{black}Treatment  & \color{black}200,000   & {\color{black}{28}}\\
      \color{black}Control    & \color{black}200,000   & {\color{black}{71}}\\
      No consent & 350,000   & 46\\
   \end{tabular}
&\hspace{20pt}
   \begin{tabular}{lcc}
\hline
                 & Size        &  Rate${}^*$\vphantom{\LARGE Y}\\
      \color{black}Grade 2 (vaccine)\hfill        & \color{black}225,000 & 25\\
      \color{black}Grades 1 and 3 (control)        & \color{black}725,000 & {\color{black}{54}}\\
      Grade 2 (no consent)\hfill     & 125,000 & {\color{black}{44}}\\
   \end{tabular}
\end{tabular}}\vspace{-4pt}
\end{center}
{${}^*$Rate is number of cases per 100,000 subjects.}
\smallskip

Describe the treatment and control groups and how they were selected 
in the two experiments.  
What does \textit{consent} mean in the context of the experiments?
Which rates demonstrate the vaccine's effectiveness in the two experiments 
(there are two rates to compare for each experiment)?  
Which experiment gave stronger evidence for the vaccine's effectiveness and why?
What does \textit{double-blind} mean?
What impact did the different sizes of the six groups have on your conclusions and why?
(20 points)


\vfill
\eject

\pagestyle{empty}
\lstset{language=R, showspaces=false, showstringspaces=false}

\href{http://www.shepherd.edu}{\includegraphics[height=1.75cm]{logo-high-res.eps}}
\vspace{-1.69cm}

{\small
\begin{tabular}{cl}
& Math 314\\
& Statistics\\
\hspace{5.28in} & 19 September 2011
\end{tabular}
}
\setlength{\baselineskip}{1.05\baselineskip}
\bigskip

\begin{center}
\textbf{\large  Exam I}
\end{center}
\medskip

1. For the $(x,y)$ data shown in the table, do the following.

\hspace{20pt} (a) Draw a scatter plot. (10 points)

\hspace{20pt} (b) Calculate the correlation coefficient. (10 points)

\hspace{20pt} (c) Find an equation for the SD line. (10 points)

\hspace{20pt} (d) Plot the SD line.  (5 points)

\rput(14.3,2){\begin{tabular}{|r|lllllll|}\hline
$x$ & 1 & 2 & 3 & 4 & 5 & 6 & 7\vphantom{\Large Y}\\[3pt]
$y$ & 7 & 6 & 5 & 4 & 3 & 1 & 2\\[1pt]\hline
\end{tabular}}
\vfill
\eject

2. (Hypothetical).  In a clinical trial, data collection usually starts at
\lq\lq baseline,\rq\rq\ when the subjects are 
recruited into the trial but before they are randomized to 
treatment or control.  Data collection continues until the end of the followup.
Two clinical trials on prevention of heart attacks report baseline data 
on weight, shown below.  In one of these trial, the randomization did not work.
Which  one and why?  (15~points)\vspace{-15pt}
\begin{center}
\begin{tabular}{cc}
\begin{tabular}{l}
\\
\\
\rput(-12pt,-3pt){(i)}
Treatment\vphantom{\Large Y}\\
Control\\[10pt]
\rput(-12pt,-3pt){(ii)}
Treatment\\
Control
\end{tabular}
&
\begin{tabular}{ccc}
Number of & Average &\\
persons & weight    & SD\\\hline
1,012   & 185 lb    & 25 lb\vphantom{\Large Y}\\
\Zc997  & 143 lb    & 26 lb\\[10pt]
\Zc995  & 166 lb    & 27 lb\\
1,017   & 163 lb    & 25 lb
\end{tabular}
\end{tabular}
\end{center}
\vspace{1.5in}

3. In 2005, men averaged about 540 on the Math SAT.  Estimate the percentage of 
men who scored over~420.  You may assume (i) that the histogram 
followed the normal curve and (ii)~the SD was about~120.  (15~points)
\vspace{2.25in}

4. For the list 1, 3, 5, 7, 9, the mean is 5 and the SD is about 2.8.

\hspace{20pt} a) What are the median, minimum, maximum, 25th percentile and the 75th percentile?
(5 points)
\vspace{1in}

\hspace{20pt} b) What are the mean and SD of the list 5, 7, 9, 11, 13?  
Explain briefly. (5 points)
\vspace{0.75in}

\hspace{20pt} c) What are the mean and SD of the list 2, 6, 10, 14, 18?  Explain briefly. (5 points)
\vfill
\eject


5. The table below shows results from two experiments designed to test the Salk vaccine in  1954.
\vspace{-5pt}
\begin{center}
{%\footnotesize
\begin{tabular}{cc}
Randomized controlled          & \\
double-blind experiment        & NFIP study\\
   \begin{tabular}{lcc}\hline
                 & Size      & Rate${}^*$\vphantom{\LARGE Y}\\
      \color{black}Treatment  & \color{black}200,000   & {\color{black}{28}}\\
      \color{black}Control    & \color{black}200,000   & {\color{black}{71}}\\
      No consent & 350,000   & 46\\
   \end{tabular}
&\hspace{20pt}
   \begin{tabular}{lcc}
\hline
                 & Size        &  Rate${}^*$\vphantom{\LARGE Y}\\
      \color{black}Grade 2 (vaccine)\hfill        & \color{black}225,000 & 25\\
      \color{black}Grades 1 and 3 (control)        & \color{black}725,000 & {\color{black}{54}}\\
      Grade 2 (no consent)\hfill     & 125,000 & {\color{black}{44}}\\
   \end{tabular}
\end{tabular}}\vspace{-4pt}
\end{center}
{${}^*$Rate is number of cases per 100,000 subjects.}
\smallskip

Describe the treatment and control groups and how they were selected 
in the two experiments.  
What does \textit{consent} mean in the context of the experiments?
Which rates demonstrate the vaccine's effectiveness in the two experiments 
(there are two rates to compare for each experiment)?  
Which experiment gave stronger evidence for the vaccine's effectiveness and why?
What does \textit{double-blind} mean?
What impact did the different sizes of the six groups have on your conclusions and why?
(20 points)


\vfill
\eject

\pagestyle{empty}
\lstset{language=R, showspaces=false, showstringspaces=false}

\href{http://www.shepherd.edu}{\includegraphics[height=1.75cm]{logo-high-res.eps}}
\vspace{-1.69cm}

{\small
\begin{tabular}{cl}
& Math 314\\
& Statistics\\
\hspace{5.28in} & 19 September 2011
\end{tabular}
}
\setlength{\baselineskip}{1.05\baselineskip}
\bigskip

\begin{center}
\textbf{\large  Exam I}
\end{center}
\medskip

1. For the $(x,y)$ data shown in the table, do the following.

\hspace{20pt} (a) Draw a scatter plot. (10 points)

\hspace{20pt} (b) Calculate the correlation coefficient. (10 points)

\hspace{20pt} (c) Find an equation for the SD line. (10 points)

\hspace{20pt} (d) Plot the SD line.  (5 points)

\rput(14.3,2){\begin{tabular}{|r|lllllll|}\hline
$x$ & 1 & 2 & 3 & 4 & 5 & 6 & 7\vphantom{\Large Y}\\[3pt]
$y$ & 6 & 7 & 5 & 4 & 3 & 1 & 2\\[1pt]\hline
\end{tabular}}

\vfill
\eject
2. (Hypothetical).  In a clinical trial, data collection usually starts at
\lq\lq baseline,\rq\rq\ when the subjects are 
recruited into the trial but before they are randomized to 
treatment or control.  Data collection continues until the end of the followup.
Two clinical trials on prevention of heart attacks report baseline data 
on weight, shown below.  In one of these trial, the randomization did not work.
Which  one and why?  (15~points)\vspace{-15pt}
\begin{center}
\begin{tabular}{cc}
\begin{tabular}{l}
\\
\\
\rput(-12pt,-3pt){(i)}
Treatment\vphantom{\Large Y}\\
Control\\[10pt]
\rput(-12pt,-3pt){(ii)}
Treatment\\
Control
\end{tabular}
&
\begin{tabular}{ccc}
Number of & Average &\\
persons & weight    & SD\\\hline
1,012   & 169 lb    & 25 lb\vphantom{\Large Y}\\
\Zc997  & 167 lb    & 26 lb\\[10pt]
\Zc995  & 188 lb    & 27 lb\\
1,017   & 148 lb    & 25 lb
\end{tabular}
\end{tabular}
\end{center}
\vspace{1.5in}

3. In 2005, women averaged about 500 on the Math SAT.  Estimate the percentage of 
women who scored over~620.  You may assume (i) that the histogram 
followed the normal curve and (ii)~the SD was about~120.  (15~points)
\vspace{2.25in}

4. For the list 1, 3, 5, 7, 9, the mean is 5 and the SD is about 2.8.

\hspace{20pt} a) What are the median, minimum, maximum, 25th percentile and the 75th percentile?
(5 points)
\vspace{1in}

\hspace{20pt} b) What are the mean and SD of the list 4, 6, 8, 10, 12?  
Explain briefly. (5 points)
\vspace{0.75in}

\hspace{20pt} c) What are the mean and SD of the list 3, 9, 15, 21, 27?  Explain briefly. (5 points)
\vfill
\eject


5. The table below shows results from two experiments designed to test the Salk vaccine in  1954.
\vspace{-5pt}
\begin{center}
{%\footnotesize
\begin{tabular}{cc}
Randomized controlled          & \\
double-blind experiment        & NFIP study\\
   \begin{tabular}{lcc}\hline
                 & Size      & Rate${}^*$\vphantom{\LARGE Y}\\
      \color{black}Treatment  & \color{black}200,000   & {\color{black}{28}}\\
      \color{black}Control    & \color{black}200,000   & {\color{black}{71}}\\
      No consent & 350,000   & 46\\
   \end{tabular}
&\hspace{20pt}
   \begin{tabular}{lcc}
\hline
                 & Size        &  Rate${}^*$\vphantom{\LARGE Y}\\
      \color{black}Grade 2 (vaccine)\hfill        & \color{black}225,000 & 25\\
      \color{black}Grades 1 and 3 (control)        & \color{black}725,000 & {\color{black}{54}}\\
      Grade 2 (no consent)\hfill     & 125,000 & {\color{black}{44}}\\
   \end{tabular}
\end{tabular}}\vspace{-4pt}
\end{center}
{${}^*$Rate is number of cases per 100,000 subjects.}
\smallskip

Describe the treatment and control groups and how they were selected 
in the two experiments.  
What does \textit{consent} mean in the context of the experiments?
Which rates demonstrate the vaccine's effectiveness in the two experiments 
(there are two rates to compare for each experiment)?  
Which experiment gave stronger evidence for the vaccine's effectiveness and why?
What does \textit{double-blind} mean?
What impact did the different sizes of the six groups have on your conclusions and why?
(20 points)

\vfill
\eject

\end{document}
